\chapter{Installing and running}\label{ch:install}
\begin{goals}
\textbf{In this chapter} you choose how to run lamina (a web address, an installed app, or the file
itself), learn what it keeps on your machine between sessions, and how to clear it.
\end{goals}

lamina is one HTML file. Every way of running it runs that file; they differ only in how it reaches you
and what the browser lets it do around it.

\section{From the web}

lamina is published at \file{gentropic.org/lamina}. Open the address in a current Chrome, Edge or
Firefox and it runs. Opening the page doesn't send your data anywhere: the files you open are read by
the page, in the browser, and the page can't make network connections (\cref{ch:sealed}).

\subsection{Installing it as an app}

Chrome and Edge offer to install the page as an app (the install icon at the end of the address bar).
Installed, lamina opens in its own window, without the browser's tabs and toolbar, and works offline:
the page keeps a copy of itself so it starts with no connection. When a newer lamina has been published,
the installed copy picks it up on a later start.

\subsection{The Microsoft Store}

On Windows, lamina is also in the Microsoft Store as \emph{GCU Lamina}. It is the same page, packaged as
an installed app. The Store version and the browser-installed version are the same lamina; choose
whichever your machine's policy allows.

\section{The file itself}\label{sec:single-file}

The empty screen offers \ui{download lamina.html}. That file is all of lamina: keep it in a folder, copy
it onto a shared drive, attach it to an e-mail, and open it by double-clicking. It runs from the disk
(the address starts \file{file://}), with no installation and no server, on a machine with no network at
all.

Two things differ when lamina runs from the disk:

\begin{itemize}
\item Browsers don't let a page on \file{file://} start background workers from inside itself, so lamina
indexes a file on the page's own thread instead. The result is the same; on a very large file the page
is less responsive while the index is built, and the status line counts the progress.
\item Chrome and Edge give a page on \file{file://} the same file-picker and save-dialog access as one on
the web, so recents and export work as usual (\cref{sec:recents}, \cref{ch:export}). Firefox has no
save dialog for pages, so an export there is downloaded instead.
\end{itemize}

\begin{tip}
The file on disk is the one to choose when a machine is locked down or air-gapped, or when you need to
show someone exactly which lamina you used. Its build stamp and its SHA-256 identify it
(\cref{sec:verify}).
\end{tip}

\section{More than one window}

\menu{File>New window} opens a second lamina, with nothing open in it. Within one window, files open in
tabs (\cref{ch:views}); a second window is for keeping two sets of tabs apart, or on two screens. The
windows don't share anything but the recents list and the index cache.

\section{Inside the lead-acid shell}

On an Android phone, lamina also runs inside lead-acid, the GCU app that hosts GCU tools with access to
the phone's files. There, \emph{Open with} lamina from the phone's file manager or a \emph{Share} to
lamina opens the file in place, read by range from the phone's storage rather than copied, and exports
go to the phone's \file{Downloads} folder, from where they can be shared. Everything else is as on a
computer.

\section{Inside another page}

A web page can embed lamina in a frame and drive it: open a file, apply a filter, read rows back. That is
for people building tools around lamina; \cref{app:host} documents it.

\section{What lamina keeps on your machine}\label{sec:local-state}

lamina remembers a few things between sessions, all in the browser's own storage for the page, on this
machine. None of it can leave it.

\begin{center}\small
\begin{tabularx}{\linewidth}{@{}lL@{}}
\toprule
\sffamily\bfseries what & \sffamily\bfseries why, and how to remove it \\
\midrule
the index of each file opened & so a file opens at once the second time (\cref{sec:index-cache}); keyed
by the file's name, size and modification time, so a changed file is indexed afresh. Holds where rows
start, not their contents. \\
the recent files & names and, where the browser allows, handles to reopen them with permission
(\cref{sec:recents}). \menu{File>Clear recents}, or turn \menu{File>Remember recent files} off. \\
each file's last views, and the last session & so a file reopens as you left it, and the empty screen can
offer the files you had open (\cref{sec:last-views}). Off with \menu{File>Remember recent files}. \\
your settings & the theme, whether distributions are drawn and log-scaled, scroll lock, the brush mode,
the width of the side panels. \\
\bottomrule
\end{tabularx}
\end{center}

The index cache and the settings go when you clear the site's data in the browser (in Chrome and Edge,
the padlock or page-info icon in the address bar, then \emph{Site settings} and \emph{Delete data}).
Each copy of lamina keeps its own: lamina from the web, the installed app, and \file{lamina.html} on
disk are different origins to the browser.
